\documentclass[11pt]{article}
\usepackage{preamble}
\setlength{\parskip}{4pt}

\begin{document}

\daybanner{AP Statistics \textperiodcentered\ Unit 3 \textperiodcentered\ Topic 3.9 \textperiodcentered\ B: 2026-12-10 \textperiodcentered\ E: 2027-01-20 \textperiodcentered\ TEACHER KEY}{Two Ways to Sample Library Visits}

\sectionbanner{CED TETHER}

{\small
\begin{itemize}[leftmargin=1.5em,itemsep=2pt,topsep=2pt]
  \item \textbf{UNC-3.N} Determine parameters of a sampling distribution for a difference in sample proportions. \textbf{[Skill 3.B]}
  \item \textbf{UNC-3.N.1} For a categorical variable, when randomly sampling from two independent populations with proportions $p_1$ and $p_2$, the sampling distribution of the difference in sample proportions has mean $\mu=p_1-p_2$.
  \item \textbf{UNC-3.N.2} If sampling without replacement and sample sizes are less than 10\% of population sizes, the difference is negligible.
  \item \textbf{UNC-3.O} Determine whether a sampling distribution for a difference of sample proportions can be described as approximately normal. \textbf{[Skill 3.C]}
  \item \textbf{UNC-3.O.1} The sampling distribution of the difference in sample proportions will have an approximate normal distribution provided sample sizes are large enough: $n_1p_1\geq10$, $n_1(1-p_1)\geq10$, $n_2p_2\geq10$, $n_2(1-p_2)\geq10$.
  \item \textbf{UNC-3.P} Interpret probabilities and parameters for a sampling distribution for a difference in proportions. \textbf{[Skill 4.B]}
  \item \textbf{UNC-3.P.1} Parameters for a sampling distribution for a difference of proportions should be interpreted using appropriate units and within the context of specific populations.
\end{itemize}
}
\textbf{Essential question:} When is the independent two-proportion sampling model warranted, and what changes for paired observations?\par
\textbf{DOK-3 skill:} 4.B \textperiodcentered\ \textbf{FRQ pattern:} {\small\texttt{independent-\allowbreak{}vs-\allowbreak{}paired-\allowbreak{}which-\allowbreak{}sampling-\allowbreak{}distribution-\allowbreak{}model-\allowbreak{}is-\allowbreak{}warranted}}\par
\textbf{DOK rationale:} Part (c) compares independent and repeated-person sampling plans, evaluates the useful and unsupported parts of competing claims, and explains how a mismatched model can misstate uncertainty.\par

\frameworkphaseheader{Do Now (first take) $\rightarrow$ Explore (video + follow-along) $\rightarrow$ Exit (finish + turn in)}{1 $\rightarrow$ 3}{45}{%
\begin{itemize}[leftmargin=1.5em,itemsep=2pt,topsep=2pt]
  \item Board slide up at the bell. Ask students to mark one sound idea and one overreach across the two claims.
  \item Begin the lesson video follow-along at minute 5.
  \item Return to the board slide at minute 33. Circulate and attach every condition to the formula feature it supports.
  \item Collect at the door. Score part (c) only, E/P/I.
\end{itemize}
}{%
\begin{itemize}[leftmargin=1.5em,itemsep=2pt,topsep=2pt]
  \item Choose whose use of the formula is warranted and cite one design detail before the lesson.
  \item Complete the follow-along about center, spread, shape, and probabilities for differences in sample proportions.
  \item Finish (a)--(c), revise the first take in the exit line, and turn in the sheet.
\end{itemize}
}{%
\begin{itemize}[leftmargin=1.5em,itemsep=2pt,topsep=2pt]
  \item What random object is selected in Plan I, and what is selected in Plan P?
  \item Which four expected counts establish the normal shape for Plan I?
  \item Does using 120 twice make the two results independent?
  \item What information about within-cardholder changes disappears if the plans are treated alike?
\end{itemize}
}{Solo teacher. Circulate ~5 pairs. Require students to separate a useful paired design from an invalid use of the independent-sample variance formula.}

\begin{callout}{\IconWarn\ WATCH FOR}{calloutred}
\begin{itemize}[leftmargin=1.5em,itemsep=2pt,topsep=2pt]
  \item Treating equal sample sizes as evidence of independence.
  \item Subtracting the two variance terms because the proportions are subtracted.
  \item Calling Plan P unusable instead of identifying the need for a paired analysis.
\end{itemize}
\end{callout}

\sectionbanner{SCORING PART (c) --- E / P / I}

\textbf{Expected elements}
\begin{itemize}[leftmargin=1.5em,itemsep=2pt,topsep=2pt]
  \item \textbf{judgment-1} (required) --- Applies the independent model to Plan I using separate random selection and passing conditions.
  \item \textbf{judgment-2} (required) --- Recognizes the useful within-person comparison in Plan P while explaining its linked responses.
  \item \textbf{judgment-3} (required) --- Explains that equal sample sizes do not imply independence and transferring the SD/tail can misstate uncertainty.
\end{itemize}

{\small\renewcommand{\arraystretch}{1.25}
\noindent\begin{tabular}{|p{0.5in}|p{5.9in}|}
\hline
\textbf{E} & Meets all three checks with correct contextual reasoning; equivalent wording is acceptable. \\ \hline
\textbf{P} & Shows at least one substantive correct check, but has an omission or error that prevents E. \\ \hline
\textbf{I} & Shows none of the three checks with substantive correct reasoning; a choice alone is insufficient. \\ \hline
\end{tabular}}

\textbf{Common mistakes}
\begin{itemize}[leftmargin=1.5em,itemsep=2pt,topsep=2pt]
  \item Treating equal sample sizes as an independence condition
  \item Subtracting the variance contributions because the proportions are subtracted
  \item Using the independent-sample normal probability for the same 120 people
  \item Calling the paired design invalid rather than changing the analysis
\end{itemize}

\clearpage
\focusbanner{TODAY'S PROBLEM --- annotated}

A library studies self-checkout before and after a redesign. Records give $p_B=0.42$ among 4{,}000 before visits and $p_A=0.54$ among 5{,}000 after visits.\par\smallskip \textbf{Plan I:} Draw an SRS of 120 before visits and a separate SRS of 150 after visits; do not track cardholders.\par \textbf{Plan P:} Measure the same 120 cardholders before and after.\par\smallskip \textbf{Sora:} Separate selection warrants Plan I's independent model; 0.0350 means a difference at most 0.01 is unusual.\par \textbf{Mateo:} Pairing controls person differences, so Plan P may help. Because both proportions use 120 people, Plan I's model applies to P too.\par\smallskip For Plan I, the four expected counts are 81, 69, 50.4, and 69.6. A normal calculator gives the lower-tail probability 0.0350.\par
\begin{center}
\textbf{Sampling plans}\par\smallskip
\begin{tabular}{|l|c|c|c|}
\hline
\textbf{Plan} & \textbf{Before} & \textbf{After} & \textbf{Repeated} \\ \hline
\textbf{I} & SRS 120 & SRS 150 & no \\ \hline
\textbf{P} & 120 & same 120 & yes \\ \hline
\end{tabular}
\end{center}
\begin{firsttakebox}{FIRST TAKE (5 min)}
Does using the same people twice seem different from choosing two separate groups? Give one reason.\par
\answer{Not graded. Sora matches the model; Mateo values a sound design but transfers the wrong model.}
\end{firsttakebox}

\textit{Rules callout ``THE VARIANCE FORMULA REQUIRES INDEPENDENT SAMPLES (keep on the page)'' is printed on the student sheet.}\par\medskip

\dokbadge{1}~\textbf{(a)} For after minus before, identify the parameter and statistic and state Plan I's sampling-distribution mean.\par
\answer{The parameter $p_A-p_B$ is the after-minus-before population difference; the statistic is $\widehat p_A-\widehat p_B$. Its mean is $0.54-0.42=0.12$.}

\dokbadge{2}~\textbf{(b)} For Plan I, verify both 10\% conditions and use the four supplied counts to check normality. Calculate the SD and interpret the supplied probability $P(\widehat p_A-\widehat p_B\leq0.01)=0.0350$.\par
\answer{Both 10\% checks pass: $150<500$ and $120<400$. All four supplied counts exceed 10. The SD is $\sqrt{0.54(0.46)/150+0.42(0.58)/120}=0.0607124$. About 3.5\% of independent sample pairs would show an after-minus-before difference at most 0.01 when the population difference is 0.12.}

\dokbadge{3}~\textbf{(c)} In 3--4 sentences, explain what Sora and Mateo each get right and wrong. Use how the samples are selected to decide where the independent model applies. Explain why transferring its SD and probability to the repeated people could mislead a reader.\par
\answer{Sora's model fits Plan I: separate SRSs and passing conditions support the independent formula. Mateo is right that repeated people can help compare within-person changes. However, their linked responses are dependent, even with equal sample sizes. Transferring the independent SD and tail probability to Plan P ignores those links and can misstate uncertainty.}

\begin{summaryexitbox}{EXIT}
One thing the video changed about my first take: \hrulefill
\end{summaryexitbox}

\end{document}
